\documentclass[UTF8,11pt,fontset=windows]{ctexart}
\usepackage[a4paper,left=24mm,right=24mm,top=23mm,bottom=24mm]{geometry}
\usepackage{amsmath,amssymb,amsthm,mathtools,bm}
\usepackage{booktabs,tabularx,array,enumitem,microtype,xcolor,needspace}
\usepackage[numbers,sort&compress]{natbib}
\usepackage[colorlinks=true,linkcolor=blue!45!black,citecolor=blue!45!black,urlcolor=blue!45!black]{hyperref}
\usepackage{fancyhdr}
\hypersetup{pdftitle={WZZT 的结构、对偶筛选与块必要性},pdfauthor={数学研究稿}}
\setlength{\headheight}{15pt}\pagestyle{fancy}\fancyhf{}
\fancyhead[L]{\small WZZT：结构与可复用证书}\fancyhead[R]{\small 2026-10-28}
\fancyfoot[C]{\thepage}
\setlength{\parskip}{0.25em}\setlength{\emergencystretch}{1.5em}
\setlist{nosep,leftmargin=2em,topsep=3pt,itemsep=2pt}
\renewcommand{\arraystretch}{1.18}\Urlmuskip=0mu plus 1mu
\numberwithin{equation}{section}
\newtheorem{definition}{定义}[section]
\newtheorem{lemma}[definition]{引理}
\newtheorem{proposition}[definition]{命题}
\newtheorem{theorem}[definition]{定理}
\newtheorem{corollary}[definition]{推论}
\theoremstyle{remark}\newtheorem{remark}[definition]{说明}
\newtheorem{example}[definition]{例}\renewcommand{\proofname}{证明}
\newcommand{\RR}{\mathbb R}\newcommand{\one}{\bm 1}
\newcommand{\diag}{\operatorname{diag}}\newcommand{\vecop}{\operatorname{vec}}
\newcommand{\pos}[1]{\left(#1\right)_{+}}
\newcommand{\nn}[1]{\left\lVert#1\right\rVert}
\newcommand{\ip}[2]{\left\langle#1,#2\right\rangle}
\newcommand{\calF}{\mathcal F}\newcommand{\calD}{\mathcal D}
\newcommand{\calM}{\mathcal M}\newcommand{\calU}{\mathcal U}
\newcommand{\sav}[1]{\left\lVert#1\right\rVert_{\mathrm{sum},1}}
\title{WZZT 的结构、对偶筛选与块必要性\\[3pt]\large 当前模型、完整推导与可拓展边界}
\author{数学研究稿}\date{2026年10月28日}
\begin{document}
\maketitle
\begin{abstract}
本文将当前 WZZT 的实际求解结构与两条可拓展方向对齐：固定支撑 L1 线性规划的安全候选筛选，以及允许其他线路补全、替换后的块必要性检查。先定义共同 R/X 树模型，推导盒约束 LP 对偶与可复用筛选界；再给出完整结构域的参考 MILP、共同三点约束及禁块条件，区分精确模型与安全外松弛。三终端算例贯通 NJ、非最优旧对偶的筛选价值和禁块证书；反例明确秩、正对偶内积及交叉分解的判断边界。本文包含源码静态核对和已有算式的精确算术复核，不包含新优化实验或加速结果。
\end{abstract}
\section{当前结构与本文要解决的问题}
\label{sec:current}
\subsection{量测、树响应与识别对象}
令终端集合为 $L=\{1,\ldots,n\}$，已观测根为 $0$。在 $T$ 个同步时刻定义
\begin{equation}
 P=[p_1,\ldots,p_T],\quad Q=[q_1,\ldots,q_T],\quad
 Y_{it}=V_0(t)^2-V_i(t)^2,\qquad P,Q,Y\in\RR^{n\times T}.
 \label{eq:data}
\end{equation}
采用负荷为正的一阶平方电压模型，内部注入为零或已计入模型：
\begin{equation}
 Y=RP+XQ+E,\qquad
 R=\sum_{C\in\calF}r_Cz_Cz_C^\top,\quad
 X=\sum_{C\in\calF}x_Cz_Cz_C^\top.
 \label{eq:model}
\end{equation}
其中 $z_C\in\{0,1\}^n$ 是非空下游集合 C 的指示向量，$r_C,x_C\ge0$。平方电压响应权重吸收系数 2，结合标幺基准才能换算为物理阻抗。支撑族 $\calF$ 必须层状：任意两个成员相交时，一方包含另一方。R 与 X 使用相同支撑族，各自权重可以不同、可以为零。

本文使用逐元素绝对值之和
\begin{equation}
 L(R,X)=\sav{Y-RP-XQ}:=\sum_{i,t}|Y_{it}-(RP+XQ)_{it}|.
 \label{eq:loss}
\end{equation}
代码目标为 $L/N$，$N=nT$，多场景时为全部标量观测数。证明与阈值必须采用同一归一化。本文把观测 P/Q 固定为回归输入，不含未知源或潜在 P/Q 的联合优化。结论首先属于式~\eqref{eq:model} 的识别域；AC 失配和量测不确定性应由误差预算或另行扩大的模型处理。

\begin{definition}[有效块与约化树]
$r_C+x_C>0$ 的块称为有效块；两通道权重均为零的辅助块不计为存在。同支撑的串联线路合并，零增量隐藏层收缩。识别对象是由根和终端确定的约化树，不区分产生相同端口核的隐藏二度串联细节。
\end{definition}

\subsection{当前实现的优化顺序}
权威实现为 \path{rnj_wzzt_core/rnj_wzzt}。当前主线为观测根平方电压差、约束 R/X 回归、RX75--RNJ 与块重采样筛选初始块、层状 L1 单块扩展，再在独立验证上用 one-SE 规则回选路径点。one-SE 指在最小验证误差的一个标准误范围内选较简单的路径模型。

\begin{table}[htbp]
\centering\small
\begin{tabularx}{\linewidth}{@{}p{0.21\linewidth}XX@{}}
\toprule
层次 & 当前优化对象 & 证书覆盖范围\\
\midrule
固定支撑重估 & 已有全部 R/X 权重的有界 L1 LP & 指定数据、支撑族及参数界\\
单次扩展 & 一个新块；全部权重共同重估 & 固定旧族下的有界单块扩展\\
完整前向路径 & 依次接受扩展；零权重重解无损时剪除非保护块 & 贪心路径，不是全树联合最优\\
初始块与验证 & 初始块固定保留；路径中回选 & 不提供替换旧分支的搜索\\
\bottomrule
\end{tabularx}
\caption{当前实现与证书对象。源码位置见附录~\ref{app:source}。}
\end{table}

当前 MILP 用 $z_i$ 表示新块成员，用 $h_{ij}=z_i z_j$、$v^R_{ij}=r_+h_{ij}$ 和 $v^X_{ij}=x_+h_{ij}$ 提升乘积。整数解处，$v^R,v^X$ 是新支路的核矩阵贡献；LP 松弛通常不对应真实支路。旧权重已共同重估，仍缺少正权重错误块的主动替换和完整替代树域审查。数值上界扩大时，前向算法从初始族重启。

\subsection{与 NJ 的差异及两条拓展方向}
任一单通道核可转换为
\begin{equation}
 d_{ij}=R_{ii}+R_{jj}-2R_{ij}
 =\sum_Cr_C|z_{C,i}-z_{C,j}|,\qquad d_{0i}=R_{ii}.
 \label{eq:kernelmetric}
\end{equation}
反向有 $R_{ij}=(d_{0i}+d_{0j}-d_{ij})/2$。完整精确的根锚定距离下，二者描述相同约化树。含噪时，WZZT 按量测拟合，NJ 按距离组合递归合并；四节点 NJ 与四点判据等价\citep[\S3, Proposition~6]{mihaescu2007}。

以仅有功的平方损失为例，若 $PP^\top$ 可逆，令 $\widehat R=YP^\top(PP^\top)^{-1}$，则
\begin{equation}
 \nn{Y-RP}_F^2=\nn{Y-\widehat RP}_F^2+
 \operatorname{tr}\!\left[(R-\widehat R)PP^\top(R-\widehat R)^\top\right].
 \label{eq:inputmetric}
\end{equation}
展开平方并用 $(Y-\widehat RP)P^\top=0$ 即可证明。直接拟合按输入确定的尺度比较核误差；保留 $\widehat R$ 和 $PP^\top$ 的后处理也能使用同一准则，因此这不是不可转移的信息优势。

\textbf{方向 A} 保留当前扩展域，用 LP 对偶减少候选求解；\textbf{方向 B} 允许其他块完整补全和替换，用有效下界判断目标块是否必要。A 的局部无增益不能替代 B 的全域必要性。

\section{固定支撑 LP 的对偶与安全候选筛选}
\label{sec:dual}
\subsection{统一特征与对偶推导}
取与源码等价的固定向量化顺序，令
\begin{equation}
 b=\vecop(Y),\quad
 a_C^P=\vecop(z_Cz_C^\top P),\quad
 a_C^Q=\vecop(z_Cz_C^\top Q).
 \label{eq:features}
\end{equation}
堆叠旧族 B 的两通道列为 $A_B$，对应权重为 $\theta$，正的有限上界为 $M_B$：
\begin{equation}
 J_B=\min_{0\le\theta\le M_B}\nn{b-A_B\theta}_1.
 \label{eq:fixed}
\end{equation}
\begin{proposition}[盒约束 L1 LP 对偶]
定义 $\calU=\{u:\nn{u}_\infty\le1\}$，则
\begin{equation}
 J_B=\max_{u\in\calU}D_B(u),\qquad
 D_B(u)=b^\top u-\sum_{k\in B}M_k\pos{a_k^\top u}.
 \label{eq:dual}
\end{equation}
求和包含每块的 P/Q 两列。
\end{proposition}
\begin{proof}
引入自由残差 e 和约束 $A_B\theta+e=b$。以 u 乘 $b-A_B\theta-e$，拉格朗日函数为 $\nn e_1+u^\top(b-A_B\theta-e)$。对 e 取下确界，$\nn u_\infty\le1$ 时结果为零，否则为 $-\infty$。再对 $\theta_k\in[0,M_k]$ 最小化，得到 $-M_k\pos{a_k^\top u}$。原 LP 总可行且最优值有限，强对偶给出等式。
\end{proof}
移除有限上界时，对偶改为 $\max b^\top u$，约束 $\nn u_\infty\le1$、$A_B^\top u\le0$。有界模型中任意盒内 u 均可通过惩罚项产生合法下界；无界模型不能省略列可行性检查。

残差非零处 $u_\ell=\operatorname{sign}(e_\ell)$，零残差处可取 $[-1,1]$ 内的值。$0<\theta_k<M_k$ 时有 $a_k^\top u=0$。u 是最优拟合损失对量测的次梯度，权重边界乘子是放宽参数域的损失敏感度；它们通常不是电价、物理电流或线路存在概率。整数 MILP 本身也没有这一 LP 影子价格解释。

\subsection{安全界、并列结构与证书复用}
\begin{theorem}[允许旧权重共同重估的筛选界]
对任意 $u\in\calU$ 和合法新块 C，
\begin{equation}
 J_{B\cup C}\ge\ell_C(u):=D_B(u)
 -M_C^P\pos{(a_C^P)^\top u}
 -M_C^Q\pos{(a_C^Q)^\top u}.
 \label{eq:screen}
\end{equation}
\end{theorem}
\begin{proof}
新增两列后，对偶仍由同一个 $\calU$ 参数化，目标恰为右端。弱对偶即得结论；推导未固定任何旧权重。
\end{proof}
保存证书库 $\{u^{(1)},\ldots,u^{(H)}\}\subset\calU$ 后，可以取
\begin{equation}
 \ell_C^{\rm bank}=\max\{0,\ell_C(u^{(1)}),\ldots,\ell_C(u^{(H)})\}
 \le J_{B\cup C}.
 \label{eq:bank}
\end{equation}
设 $U_{\rm best}$ 是同一扩展域内的合法可行目标。严格大于上界时可排除 C；相等时若只求一个最优解可以不再求它，但若保留全部并列或容差内解释，就须保留。下界不足只表示未排除。

\begin{corollary}[无增益与零权重不同]
若 u 是旧问题最优对偶，且新增两列的内积均不大于零，则 $J_{B\cup C}=J_B$。若两内积严格为负且上界为正，则两个新权重在全部扩展最优解中均为零。
\end{corollary}
\begin{proof}
前一结论由筛选界和新权重取零的可行上界夹逼得到。后一结论中，新列的上界乘子可取零，下界乘子为 $-a_k^\top u>0$；任意原始/对偶最优解的互补条件迫使相应权重为零。
\end{proof}
结论只针对固定 B 的本轮扩展；加入其他块后应重新计算，不能永久删除一个物理块。

\subsection{不能替代安全界的两个判据}
\begin{example}[正内积仍可无增益]
旧模型无变量，$b=(1,0)^\top$，最优值为 1。新增 $a=(1,1)^\top$、$0\le t\le1$ 后，目标 $|1-t|+|t|=1$。旧最优对偶可取 $(1,1)^\top$，新列内积为 2，仍无增益。取 $(1,-1)^\top$ 可证无增益，但正的新权重仍可属于并列最优解。
\end{example}
\begin{example}[秩不增加仍可改善]
一次输入 $p=(1,1)^\top$，旧族为两个 singleton，权重均在 $[0,1]$，量测为 $(2,2)^\top$。旧特征张成 $\RR^2$，可行预测却只在 $[0,1]^2$，最小损失为 2。新增共同块列 $(2,2)^\top$ 后，权重取 1 即精确拟合。秩未描述非负和上界所确定的预测集。
\end{example}
经典 Gap Safe 的球形区域利用 Lasso 平方损失的对偶曲率\citep[\S1--3]{fercoq2015}。这里的对偶分段线性且可能不唯一，不能直接移植半径公式。当前最优对偶也未必给最强候选界，第~\ref{sec:example} 节给出严格例子。

\subsection{块定价与数值证书}
将 u 还原为 $U\in\RR^{n\times T}$，则
\begin{equation}
 (a_C^P)^\top u=\sum_t(z_C^\top p_t)(z_C^\top u_t)
 =z_C^\top H_P(u)z_C,\quad
 H_P(u)=\frac{PU^\top+UP^\top}{2}.
 \label{eq:quadratic}
\end{equation}
Q 通道相同。每个证书可压缩为两个 $n\times n$ 矩阵；若候选为完整子块之并，设子块指示矩阵为 W，还可压缩为 $W^\top H_PW$。这均为精确恒等式。

对候选区域内最大内积或惩罚项，可复用 $h_{ij}=z_i z_j$ 的整数提升，省去逐时绝对残差变量。最大化的\textbf{全局上界}才能证明没有遗漏的大内积；可行候选仅提供下界。定价仍有组合代价，尚无总时间收益结论。

数值证书应先使 u 满足盒约束，再重算 $D_B$ 和全部惩罚。采用内积区间时，取 $b^\top u$ 下端和 $a_k^\top u$ 上端得到保守下界。不能用未证明相等的原始目标替代 $D_B$。参数界变化后须按新界重算。平均损失的对偶盒为 $|u_\ell|\le1/N$；追加量测时，旧平均损失证书须乘 $N_{\rm old}/N_{\rm new}$ 并对新行补零，不能原值直接复用。

\section{完整结构域中的块必要性}
\label{sec:necessity}
\subsection{目标域与上下界判定}
令 $\calD$ 为 L 的全部非空子集。给定相同、预先声明的支路上界 $M_R,M_X$，定义 $\calM_{\rm atom}$ 为式~\eqref{eq:model} 中所有层状族及其非负盒内权重产生的共同核。它允许所有遗漏块补全，并允许替换未经确认的旧块，不限于当前贪心路径前缀。
\begin{equation}
 J_{\neg C}=\min_{(R,X)\in\calM_{\rm atom},\ C\text{ 非有效块}}L(R,X).
 \label{eq:omit}
\end{equation}
有限支撑族和有限参数盒保证最小值可以取得。设 $L_{\neg C}\le J_{\neg C}\le U_{\neg C}$，且已有合法模型满足 $J_{\rm ref}\le\varepsilon$，则
\begin{equation}
 \begin{array}{ll}
 L_{\neg C}>\varepsilon &\Rightarrow C\text{ 在预算内必要},\\
 U_{\neg C}\le\varepsilon &\Rightarrow\text{存在不含 C 的合法解释},\\
 L_{\neg C}\le\varepsilon<U_{\neg C}&\Rightarrow\text{未决}.
 \end{array}
 \label{eq:classification}
\end{equation}
尚无模型满足预算时，应先报告模型失配或可行性未明，避免空集上的“必要性”。这是指定域内的确定性结论。若改用近最优性，$L_{\neg C}>J_{\rm ref}+\delta$ 可保守排除全部距全域最优值不超过 $\delta$ 的省略解释；其含义应与噪声预算分开。

\subsection{与支路盒完全对齐的有限参考 MILP}
对每个 $S\in\calD$ 引入 $s_S\in\{0,1\}$，置
\begin{equation}
 \begin{gathered}
 R=\sum_{S\in\calD}r_Sz_Sz_S^\top,\qquad
 X=\sum_{S\in\calD}x_Sz_Sz_S^\top,\\
 0\le r_S\le M_Rs_S,\quad0\le x_S\le M_Xs_S,\\
 s_S+s_D\le1\quad(S,D\text{ 交叉}),\qquad s_C=0.
 \end{gathered}
 \label{eq:dictionarymip}
\end{equation}
交叉指非空相交且互不包含。加上量测绝对残差约束后，这是与原支路盒一致的完整禁块 MILP。零权重块可以删除，故 $s_C=0$ 精确覆盖 C 非有效的全部模型。全部非空子集是数学参考域，不是截断白名单；指数规模限制其直接用途。层状族最多含 $2n-1$ 个不同非空成员，无须另设更小上限来制造分离。

\subsection{共同三点条件为何完整}
\begin{lemma}[单核层状表示]
对称 K 可表示为非负层状原子之和，当且仅当其元素非负、$K_{ii}\ge K_{ij}$，且对所有不同的 $i,j,k$ 有
\begin{equation}
 K_{ij}\ge\min\{K_{ik},K_{jk}\}.
 \label{eq:ultra}
\end{equation}
\end{lemma}
\begin{proof}
必要性来自共同祖先关系：三个非对角共享深度的最小值至少重复两次，终端自身深度不小于共享深度。
反之，对每个 $t>0$，在 $I_t=\{i:K_{ii}\ge t\}$ 上定义 $i\sim_tj$ 当且仅当 $K_{ij}\ge t$。式~\eqref{eq:ultra} 保证传递性，得到随 t 增大而细分的等价类，所有类形成层状族。有限矩阵只有有限个元素值；以每个类持续存在的 t 区间长度为权重，恒等式
\[
 K_{ij}=\int_0^\infty\one\{i,j\text{ 在同一 }t\text{ 类中}\}\,dt
\]
给出有限层状分解。对角由单元素类的持续区间恢复。
\end{proof}
树形协方差的相同表示已有文献基础\citep[Theorem~2, \S2.3]{bravo2009}。这里讨论共享路径矩阵，允许不同对角深度，并未要求叶间距离来自等深树。

\begin{theorem}[共同 R/X 层状结构]
R、X 有共同非负层状分解，当且仅当二者分别对称、非负、有序，并且每个三元组共用以下三个情形中的一个：
\begin{equation}
 \begin{array}{c|c}
 R_{ij}\ge R_{ik}=R_{jk}&X_{ij}\ge X_{ik}=X_{jk}\\
 R_{ik}\ge R_{ij}=R_{jk}&X_{ik}\ge X_{ij}=X_{jk}\\
 R_{jk}\ge R_{ij}=R_{ik}&X_{jk}\ge X_{ij}=X_{ik}
 \end{array}
 \label{eq:commontriple}
\end{equation}
\end{theorem}
\begin{proof}
共同层状树决定三终端祖先关系；各通道非负增量只能保留或压平它，因此满足共同选择。
反之，分别取引理保证的 R、X 有效层状分解。若 R 中有效块 A 与 X 中有效块 B 交叉，取 $i\in A\setminus B$、$j\in A\cap B$、$k\in B\setminus A$。A 的正增量使 $R_{ij}>R_{ik}=R_{jk}$，B 的正增量使 $X_{jk}>X_{ij}=X_{ik}$，要求两个不同的严格情形，与式~\eqref{eq:commontriple} 矛盾。故两个有效族并集层状；对各自缺少的块补零权重即可。
\end{proof}
分别约束 R、X 为树核还不够。例如 $R=I+J_{12}$、$X=I+J_{13}$ 各自合法，却要求交叉的共同块。共用分支选择会排除它们。这里及下文 $J_S=z_Sz_S^\top$。

\subsection{三点 MILP 与参数域对齐}
设有有效核矩阵界 $0\le R_{ij}\le B_R$、$0\le X_{ij}\le B_X$。每个三元组引入三个二元变量 $\rho^{(1:3)}$，其和为 1。第一种情形在 R 上例如写成
\begin{equation}
 \begin{aligned}
 R_{ik}-R_{ij}&\le B_R(1-\rho^{(1)}),\\
 R_{ik}-R_{jk}&\le B_R(1-\rho^{(1)}),\\
 R_{jk}-R_{ik}&\le B_R(1-\rho^{(1)}).
 \end{aligned}
 \label{eq:triplemip}
\end{equation}
另两种情形循环置换；X 使用同一组 $\rho$ 和自己的 $B_X$。加上量测残差后，得到 $O(n^3)$ 个分支选择的完整树\emph{结构}模型。单矩阵三点 MILP 已见\citet[Table~1]{bravo2009}；此处共同使用选择变量，并保留原始 P/Q 预测方程。

\textbf{结构完整不等于参数盒相同。} 沿终端的不同支撑链最多有 n 个成员，故每支路界推出 $R_{ij}\le nM_R$、$X_{ij}\le nM_X$。反向不成立：只满足核界时，恢复出的某条支路增量仍可能超过 $M_R$ 或 $M_X$。所以，用这些外界的三点模型包含 $\calM_{\rm atom}$，其禁块最优值可作安全下界；整数解仍须恢复共同树的增量并检查原支路盒，才能作上界。式~\eqref{eq:dictionarymip} 是小规模精确参照。

若加预算截断 $L\le\varepsilon$ 且 P/Q 均非负，则每项预测非负，$R_{ij}p_{jt}\le Y_{it}+\varepsilon$；任一 $p_{jt}>0$ 可导出核界，Q 同理。这只是预算内的有效界。存在正负注入抵消时推导失效；不能证明有限界时，应保留参数域未覆盖的结论，而不是任意截断。

\Needspace{12\baselineskip}
\subsection{禁块条件及其线性化}
令 $K=R+X$，补入根行列 $K_{0i}=K_{i0}=0$。共同族的零总增量块收缩后，K 的有效族就是两通道有效族的并集。
\begin{theorem}[列分层与块不存在]
对任意非空终端集合 C，固定任意 $j\in C$，在合法共同树核上有
\begin{equation}
 C\text{ 有效}\ \Longleftrightarrow\
 \min_{i\in C}K_{ij}>\max_{k\in(\{0\}\cup L)\setminus C}K_{kj}.
 \label{eq:clade}
\end{equation}
等价地，
\begin{equation}
 C\text{ 无效}\ \Longleftrightarrow\
 \exists i\in C,\ k\in(\{0\}\cup L)\setminus C:
 K_{ij}\le K_{kj}.
 \label{eq:omission}
\end{equation}
\end{theorem}
\begin{proof}
C 的入口总增量若为 $w_C=r_C+x_C>0$，块内 i 与 j 的共同路径至少到入口之后，块外 k 与 j 的共同路径最多到入口之前，间隔至少为 $w_C$。
反之，在严格间隔中选阈值 t，使块内列值均大于 t，块外均小于 t。K 的 t 等价类包含 j 且恰为 C；它在正长度阈值区间持续存在，因此有效。根参与比较，使全集共同干线和 singleton 的零增量边界也得到覆盖。否定严格不等式即得式~\eqref{eq:omission}。
\end{proof}
令 $B_K=B_R+B_X$，对见证对引入 $\beta_{ik}\in\{0,1\}$：
\begin{equation}
 \sum_{i\in C,\ k\notin C}\beta_{ik}=1,\qquad
 K_{ij}-K_{kj}\le B_K(1-\beta_{ik}).
 \label{eq:omitmip}
\end{equation}
k 包含根，其核值为已知零。共同三点模型加此式覆盖完整省略域，不锁定贪心旧族。若只保留有序非负矩阵约束，则是连续外松弛：各见证分支下界取\emph{最小}才覆盖析取并集；分别对不同锚点取得完整域下界后，可以取\emph{最大}加强。未检查的见证分支不能直接略去。

\section{对偶如何连接当前扩展与全域必要性}
\label{sec:bridge}
\begin{proposition}[不依赖权重上界的完整省略域证书]
设 $u\in\calU$，对所有非空 $S\ne C$ 均有
\begin{equation}
 (a_S^P)^\top u\le0,\qquad(a_S^Q)^\top u\le0.
 \label{eq:globaldual}
\end{equation}
则对全部非负权重、合法共同树、C 无效的模型，有
\begin{equation}
 J_{\neg C}\ge b^\top u.
 \label{eq:globalbound}
\end{equation}
\end{proposition}
\begin{proof}
删除零权重 C 后，每个模型只使用 $S\ne C$ 的原子。非负权重给出 $u^\top A\theta\le0$，从而
\[
 \nn{b-A\theta}_1\ge u^\top(b-A\theta)\ge b^\top u.
\]
对全部合法省略模型取最小即可。该条件还覆盖允许交叉原子的更大模型，因此是充分证书而非必要证书。
\end{proof}
可分别求解
\begin{equation}
 \max_{z\in\{0,1\}^n,\ z\ne0,\ z\ne z_C}z^\top H_P(u)z
 \label{eq:globalprice}
\end{equation}
和 Q 通道验证此条件。\textbf{此时不能加入未经确认的旧块相容限制}。它复用 $zz^\top$ 提升，却没有权重乘积和逐时残差；需用最大化的有效上界确认没有违反列。条件失败只说明该充分证书不足。

有限盒下可允许正内积。令
\[
 W_S(u)=M_R\pos{(a_S^P)^\top u}+M_X\pos{(a_S^Q)^\top u}.
\]
若 $\Gamma(u)$ 是所有合法省略族上 $\sum_SW_S(u)$ 的\emph{上界}，则 $J_{\neg C}\ge b^\top u-\Gamma(u)$。这连接候选筛选与全域补全，但计算 $\Gamma$ 本身可能困难；只计少量已见候选不保证得到上界。

\subsection{连续证书可能必然不足}
\begin{example}[交叉分解使禁块连续松弛为零]
取四终端、$P=I,Q=0,Y=R_\star$，目标 $C=\{123\}$：
\begin{equation}
 R_\star=2I+J_{1234}+J_{123}
 =\diag(1,1,1,2)+J_{1234}+J_{12}+J_{13}+J_{23}.
 \label{eq:crossing}
\end{equation}
右端没有 C，权重都非负，但二元块交叉，不是合法树。故允许交叉的禁块连续问题最优值为零，其任何对偶都不能给正下界。

合法省略树却必须有正损失。固定 $j=1$，对 $i=2,3$、$k=4$，真实列值相差 $2-1=1$，对称核总 L1 修改至少为 2；$i=1,k=4$ 的间隔为 3，修改代价至少为 3；根见证更大。因此 $J_{\neg C}\ge2$。保留共同块和 $\{12\}$、用 singleton 匹配对角，可达到损失 4，故
\begin{equation}
 2\le J_{\neg C}\le4.
 \label{eq:fourbounds}
\end{equation}
这里仅声称手推上下界，未声称禁块最优值已精确求出。
\end{example}
此例说明全原子连续松弛与列分层松弛可以互补。若缺口来自放松树结构，增加同一松弛的对偶数量也无法补上它。

\section{贯通算例：同一量测下的 NJ、筛选与必要性}
\label{sec:example}
取 $n=T=3,Q=0$：
\begin{equation}
 \begin{gathered}
 P=\diag(1,10,1),\quad
 Y=\begin{bmatrix}10&20&3\\2&100&1\\3&10&10\end{bmatrix},\\
 R_\star=\begin{bmatrix}10&2&1\\2&10&1\\1&1&10\end{bmatrix}
 =\diag(8,8,9)+J_{123}+J_{12}.
 \end{gathered}
 \label{eq:example}
\end{equation}
误差只有 $E_{13}=E_{31}=2$。统一缩小响应和量测可放入正常平方电压降范围，比较值相应缩放。误差位置仅用于解释构造，不提供给算法。下文有限盒例取共同上界 20，保证所有构造解均在域内；这不是生产默认参数的运行结果。

\subsection{输入权重使两个准则分离}
无约束最小二乘给出
\[
 \widehat R=YP^{-1}=\begin{bmatrix}10&2&3\\2&10&1\\3&1&10\end{bmatrix}.
\]
它对称、非负、有序，却不满足三点树关系。根锚定距离满足 $d_{0i}=10$、$d_{12}=16,d_{13}=14,d_{23}=18$。四点和为 $26,24,28$，NJ 选择错误的 $13\mid02$。

直接 L1 拟合则为
\begin{equation}
 \begin{aligned}
 L(R)={}&|R_{11}-10|+10|R_{22}-10|+|R_{33}-10|\\
 &+11|R_{12}-2|+2|R_{13}-3|+11|R_{23}-1|.
 \end{aligned}
 \label{eq:exampleloss}
\end{equation}
三个非对角权重 $11,2,11$ 来自同一批量测的输入幅值。

\begin{table}[htbp]
\centering
\begin{tabular}{lcc}
\toprule
完整结构类型 & 最小总 L1 & 最小平方损失\\
\midrule
共同块加 $\{12\}$ & $4$ & $808/103$\\
共同块加 $\{13\}$ & $11$ & $101/2$\\
共同块加 $\{23\}$ 或星形 & $13$ & $3737/68$\\
\bottomrule
\end{tabular}
\caption{在三种分支与星形的完整结构域上比较；对角均可匹配。}
\label{tab:examplecost}
\end{table}


$\{12\}$ 类型要求 $R_{13}=R_{23}=c\le R_{12}$。L1 取 $c=1,R_{12}=2$ 唯一达到 4，因为共享值 1 的权重 11 大于共享值 3 的权重 2。$\{13\}$ 类型要求 $R_{12}=R_{23}=c$，由 $11|c-2|+11|c-1|\ge11$ 得下界，并可达到。

对 $\{23\}$ 类型，写 $R_{12}=R_{13}=c$、$R_{23}=a\ge c$。用 $|a-1|\ge a-1\ge c-1$，分别在 $c\le2$ 和 $c\ge2$ 上展开另外两个绝对值的线性下界，得到
\[
 11|c-2|+2|c-3|+11|a-1|
 \ge\begin{cases}17-2c\ge13,&c\le2,\\20c-27\ge13,&c\ge2.\end{cases}
\]
$a=c=2$ 达到 13，故最优解落在星形边界。各类型的独立终端权重均可非负，因此上述下界可行。

平方损失的非对角权重为 $101,2,101$。$\{12\}$ 类型取两项加权均值 $c=(2\cdot3+101\cdot1)/103=107/103$，高值 $R_{12}=2$，得 $808/103$；$\{13\}$ 类型的共享值为 $3/2$，得 $101/2$。对 $\{23\}$ 类型，无约束最优为 $c=208/103,a=1$，违反 $a\ge c$；严格凸性使约束最优落在 $a=c$，加权均值为 $103/68$，得 $3737/68$。所以优势不依赖 L1 独有的鲁棒性，而在于共同树约束与量测权重的联合使用。

若每轮精确新增一个块，以三个 singleton 起步，L1 首轮 $\{123\},\{12\},\{13\},\{23\}$ 的目标为 $13,17,33,28$，分别由单列拟合或共同值的加权中位数取得。先选全集，再选 $\{12\}$ 达到 4，因此不依赖更全面的全树搜索。

改取 $P=\diag(1,s,1)$ 且 $Y=\widehat RP$，对 $s>1$，L1 两个主要类型的最优值为 $4$ 与 $s+1$，星形为 $s+3$，故严格选真分支当且仅当 $s>3$。平方损失令 $W=s^2+1$，两个主要值为 $8W/(W+2)$ 与 $W/2$，其差
\[
 \frac W2-\frac{8W}{W+2}=\frac{W(W-14)}{2(W+2)}
\]
给出阈值 $s>\sqrt{13}$；另一个错误类型仍落在星形边界，其值不小于错误 $\{13\}$ 类型。$s=1$ 时，两种损失都选错。若 NJ 前先执行同一量测加权树拟合，差异也会消失。

\subsection{当前最优对偶与旧证书的严格比较}
令 B 含三个 singleton 与全集。其最优核对角为 10、非对角均为 2，$J_B=13$。取
\begin{equation}
 U_0=\begin{bmatrix}0&9/10&1\\0&0&-1\\1&-1&0\end{bmatrix},\qquad
 U_1=\begin{bmatrix}0&1&0\\1&0&-1\\0&-1&0\end{bmatrix}.
 \label{eq:duals}
\end{equation}
两者均盒可行，对 B 的列内积均为零，故 $D_B(u_0)=13,D_B(u_1)=11$。候选分数为
\begin{equation}
 \begin{array}{c|rrr}
 &a_{12}^\top u&a_{13}^\top u&a_{23}^\top u\\\hline
 u_0&9&2&-11\\
 u_1&11&0&-11
 \end{array}
 \label{eq:scores}
\end{equation}
上界为 20 时，$u_0$ 的筛选界与零取最大后为 $0,0,13$，$u_1$ 则给 $\{13\}$ 下界 11。得到损失 4 的合法候选后，证书库可排除另外两个候选。

旧最优核的非零残差迫使所有旧最优对偶满足 $U_{13}=U_{31}=1$，正且未触界的 singleton 又迫使对角乘子为零。因此任一旧最优对偶都给 $a_{13}^\top u=2$，仅在当前最优对偶面挑选也无法改善这个界。$U_1$ 目标稍差却更有用，说明应回收不同已求解 LP 的证书。本文手工构造它以证明这一点，未声称取得它无需计算。

\subsection{同一个对偶得到完整禁块值}
目标 $C=\{12\}$ 的所有非目标原子只有三个 singleton、全集、$\{13\}$ 和 $\{23\}$。式~\eqref{eq:scores} 和旧列的零内积表明 $u_1$ 满足式~\eqref{eq:globaldual}，故 $J_{\neg C}\ge11$。合法 $\{13\}$ 树达到 11，因而
\begin{equation}
 J_{\neg\{12\}}=11,\qquad J_{\rm ref}=4.
 \label{eq:necessityexample}
\end{equation}
对 $4\le\varepsilon<11$，该块在完整非负共同树模型中必要；这份证书不需要有限权重上界。若用平均损失，两个值为 $4/9,11/9$，预算也应除以 9。

\section{合理可拓展性与验收边界}
\label{sec:extension}
\subsection{中间问题的作用域}
可行整数解可作同域上界；固定结构 LP 原始解可作重估起点；按式~\eqref{eq:dual} 重算的对偶可跨候选提供下界；结构提升可用于定价或全原子证书检查。当前 \texttt{dual\_bound} 只是 MIP 标量界，不是逐量测对偶 u；需导出并校验连续 LP 证书后才能实现上述复用。

若完整禁块域被区域 $\calM_q$ 覆盖，且每个区域有下界 $L_q$，则
\begin{equation}
 J_{\neg C}\ge\min_qL_q.
 \label{eq:coverage}
\end{equation}
同一完整域上的独立下界则取最大。证书至少保存数据及归一化、根/终端标签、参数界、物理先验、结构区域、合法上界、有效下界和停止原因。不同前向旧族的节点界不能不经覆盖证明就合成全树下界。

当前上界为 4、节点下界为 6 时，求最优可以剪枝；若预算为 10，该节点仍可能有可解释模型。必要性审查须保留全部下界不超过预算且未被其他证据排除的区域。逐步补入共同三点关系时，每一步仍是完整域的外松弛，下界有效；候选上界须通过全部树关系和原参数界检查。

\subsection{主要代价和失败条件}
\textbf{筛选收益取决于界的质量和复用成本。} 大上界会放大筛选惩罚，使下界退化为零；证书库和定价也有开销。加速应比较同数据、同参数域和同停止精度下的总时间、调用数与有效 gap，而非只报告筛除数量。原完整 LP 的逻辑推论再加一遍，一般不会提高理论松弛界。

\textbf{完整审查大于单块扩展。} 共同三点描述没有消除整数复杂度，全集参考模型具有指数列数。当前结果支持先做小规模参照和便宜下界，再对未决块进行完整审查，尚无大规模可用性结论。

\textbf{弱块可能确实不必要。} 取 $P=I,Q=0$，
\[
 R_a=I+3J_{123}+aJ_{12},\qquad
 R'_0=I+3J_{123}+a\diag(1,1,0),\quad 0<a<1.
\]
后者不含 $\{12\}$，且 $\sav{R_a-R'_0}=2a$。故无噪声参考损失为零时，$J_{\neg\{12\}}\le2a$；固定正预算下，足够弱的块无法确认为必要。这是允许的替代解释，不应以正权重下限排除。

\textbf{硬初始块可能排除真结构。} 锁定 $\{13\}$ 会使 $\{12\}$ 永远不可新增。必要性审查只能锁定明确物理先验或已在同域、同预算下确认的块。已确认块可加入块外响应相等等式，因为不移除预算内合法解释；数据、预算或参数域改变后须重核证明链。

\subsection{最小拓展顺序与当前结论}
第一步保留原前向求解器，为固定结构 LP 增加可核查对偶、原始/对偶目标和作用域标识，先旁路筛选而不改变候选合法域。第二步建立独立必要性审查接口，先组合全列对偶证书和多锚点 LP，再对未决块投入完整禁块审查，始终保存合法替代上界。这服务于接受、暂缓或撤销块，而非事后概率打分。

最低验收包括：三终端精确参照、四终端交叉松弛反例、正内积与秩不变边界，以及有界/无界、R/X 不相容、弱块和错误初始块。完整小规模参照中，每个下界不得超过对应最优值，每次排除不得删掉需要保留的并列结构。通过后再评价效率。

\begin{table}[htbp]
\centering\small
\begin{tabularx}{\linewidth}{@{}p{0.2\linewidth}XX@{}}
\toprule
内容 & 本稿已得到 & 尚未得到\\
\midrule
对偶筛选 & 对偶、允许重估的界、证书库和反例 & 实现、筛选频率及实际加速\\
完整域必要性 & 精确参考模型、共同三点与禁块条件、参数域对齐 & 大规模求解和松弛紧度实验\\
两方向连接 & 禁块值 11 的证书；四终端松弛反例 & 一般网络中的获取成本和覆盖率\\
\bottomrule
\end{tabularx}
\caption{数学结果、现有实现和待验证拓展分开报告。}
\end{table}


\appendix
\section{源码核对与交付复核范围}
\label{app:source}
核对日期为 2026-10-28，基准提交为 \texttt{10b17c16}。工作区另有既有修改；本稿只读取相关模型文件，不声称整个工作区等于该提交。所用源文件 SHA-256 保存在 \path{verification/source_snapshot.json}。

\begin{table}[htbp]
\centering\small
\begin{tabularx}{\linewidth}{@{}p{0.42\linewidth}X@{}}
\toprule
文件或位置 & 本稿使用的事实\\
\midrule
\texttt{preprocessing.py}:8--25 & 观测根平方电压差，根表误差仍进入目标\\
\texttt{laminar\_l1\_milp.py}:574--591 & 固定块的 P/Q 特征\\
同文件 712--744 & 权重盒、平均 L1 目标、固定族 LP\\
同文件 747--810 & 乘积提升、层状关系和去重\\
同文件 1187--1198 & 初始块保留、单块扩展、验证回选\\
同文件 1329--1375 & 单步界的适用域；触界扩大并重启\\
同文件 40--80、887--905 & 返回权重、目标与 MIP 标量诊断，未导出逐行对偶库\\
\texttt{pipeline.py}:332--337 & 初始族和共同权重上界传入扩展器\\
\bottomrule
\end{tabularx}
\end{table}

除 \texttt{pipeline.py} 外，表中文件位于 \path{rnj_wzzt_core/rnj_wzzt/estimation/}；\texttt{pipeline.py} 位于包根目录。生产基线、AC 实验及项目优化程序均未为本稿运行。

本目录的 \path{verify_derivations.py} 仅用有理数复核构造解、目标、对偶内积、距离和及反例恒等式，不调用优化求解器。它不替代一般定理证明，也不是算法性能实验。引文依据是 \citet{bravo2009} 的 \S2 和 Table~1、\citet{fercoq2015} 的 \S1--3，以及 \citet{mihaescu2007} 的 \S1--3，仅使用实际阅读部分，不声称综述穷尽相关工作。

\begingroup\small\raggedright\setlength{\bibsep}{2pt}
\bibliographystyle{unsrtnat}\bibliography{references}
\endgroup
\end{document}
